\documentclass[12pt,letterpaper]{article}

%Packages
\usepackage{amsmath}
\usepackage{amsthm}
\usepackage{amsfonts}
\usepackage{amssymb}
\usepackage{amscd}
\usepackage{dsfont}
\usepackage{enumerate}
\usepackage{fancyhdr}
\usepackage{mathrsfs}
\usepackage{bbm}
\usepackage{framed}
\usepackage{mdframed}
\usepackage{cancel}
\usepackage{float}
\usepackage{mathtools}
\usepackage{hyperref}
%\usepackage[]{mcode}
\usepackage{graphicx}
\usepackage{tikz}
\usetikzlibrary{automata,arrows}

%Page formatting
\usepackage[letterpaper,voffset=-.5in,bmargin=3cm,footskip=1cm]{geometry}
\setlength{\parindent}{0.0in}
\setlength{\parskip}{0.1in}
\allowdisplaybreaks
\headheight 15pt
\headsep 10pt

% Common Commands
\newcommand\N{\mathbb N}
\newcommand\Z{\mathbb Z}
\newcommand\R{\mathbb R}
\newcommand\Q{\mathbb Q}
\newcommand\lcm{\operatorname{lcm}}
\newcommand\setbuilder[2]{\ensuremath{\left\{#1\;\middle|\;#2\right\}}}
\newcommand\E{\operatorname{E}}
\newcommand\V{\operatorname{V}}
\newcommand\Pow{\ensuremath{\operatorname{\mathcal{P}}}}

\DeclarePairedDelimiter\ceil{\lceil}{\rceil}
\DeclarePairedDelimiter\floor{\lfloor}{\rfloor}

% 22 style
\newcommand\hint[1]{\textbf{Hint}: #1}
\newcommand\note[1]{\textbf{Note}: #1}
\newenvironment{22enumerate}{\begin{enumerate}[a.]\itemsep0em}{\end{enumerate}}
\newenvironment{22itemize}{\begin{itemize}\itemsep0em}{\end{itemize}}
\fancypagestyle{firstpagestyle} {
  \renewcommand{\headrulewidth}{0pt}%
  \lhead{\textbf{CSCI 0220}}%
  \chead{\textbf{Discrete Structures and Probability}}%
  \rhead{S. Haddadan}%
}



\pagestyle{fancyplain}

\newcommand\EX{\mathds{E}}
\newcommand\PR{\mathds{P}}
\newcommand\zlam{Z_\lambda}
\DeclareMathOperator*{\argmin}{arg\,min}

%%%%%%%%%%%%%%%% COMPILE WITH \soltrue or \solfalse %%%%%%%%%%%%%%%%%%%%%%%%%%%%%
\newif\ifsol
\soltrue  % SOLUTIONS
%\solfalse % NO SOLUTIONS
%%%%%%%%%%%%%%%%%%%%%%%%%%%%%%%%%%%%%%%%%%%%%%%%%%%%%%%%%%%%%%

%%%%%%%%%%%%%%%% solution help %%%%%%%%%%%%%%%%%%%%%
\newcommand{\solu}[2]{ \begin{mdframed} \ifsol #2 \else \vspace{#1} \fi \end{mdframed} }
\newcommand{\solt}{ \ifsol (T) \fi }
\newcommand{\solf}{ \ifsol (F) \fi }
\newcommand{\solm}[1]{\ifsol \textit{(#1)} \fi}


\begin{document}
  \thispagestyle{firstpagestyle}
  \begin{center}
    {\large \textbf{Recitation 10}}
    
    {\large Graph Theory}

  \end{center}

\section*{Definitions}

\begin{itemize}
    \item Let $G=(V, E)$ and $G'=(V', E')$. $G$ and $G'$ are \textbf{isomorphic} if there is a bijection $f: V \longrightarrow V'$ such that for every pair of vertices $\{x, y\} \in E$ if and only if $\{f(x), f(y)\} \in E'$. In other words, two graphs are isomorphic if we can find a consistent relabeling of the edges and vertices of the two graphs.
    \item A \textbf{walk} from $v_o$ to $v_1$ in an undirected graph $G$ is a sequence of alternating vertices and edges that starts at $v_0$ and ends with $v_1$.
    \item An \textbf{open walk} is a walk in which the first and last vertices are not the same.
    \item A \textbf{closed walk} is a walk in which the first and last vertices are the same.
    \item A \textbf{trail} is an open walk in which no edge occurs more than once.
    \item A \textbf{circuit} is a closed walk in which no edge occurs more than once.
    \item A \textbf{path} is a trail in which no vertex occurs more than once.
    \item A \textbf{cycle} is a circuit in which no vertex occurs more than once, except the first and last vertices which are the same.
    \item Two vertices are \textbf{connected} if there exists a path between them (or if they are the same vertex).
    \item An \textbf{Euler trail} is a trail that visits every edge exactly once.
    \item An \textbf{Euler circuit} is a circuit that contains every edge exactly once.
    \item A \textbf{Hamiltonion circuit} in an undirected graph is a circuit that includes every vertex (except first/last vertex) exactly once.

    \item A \textbf{Hamiltonion path} in an undirected graph is a path that includes every vertex in the graph.
    \item A \textbf{Bipartite graph} is a graph whose vertex set can be partitioned to two parts, and no edge connects a pair of vertices belonging to the same part. Any graph that is a subgraph  of a complete bipartite graph is bipartite.  
\end{itemize}
\newpage
\textbf{Task 1}

Consider the following road map:
\begin{center}\includegraphics*[scale=.4]{euler_graph.png}\end{center}
\begin{enumerate}[a.]

    \item An explorer wants to explore all the routes between a number of cities. Can a tour be found which traverses each route only once? Particularly, find a path which starts at A, goes along each road exactly once, and ends back at A.
    
    In other words, find an Euler trail for the graph!
    \begin{mdframed}\vspace{1 cm}\end{mdframed}
    
    \item A traveler wants to visit a number of cities. Can a path be found which visits each city only once? Particularly, find a tour which starts at A, goes to each city exactly once, and ends back at A.
    
    In other words, find an Hamiltonian cycle for the graph!
    \begin{mdframed}\vspace{1 cm}\end{mdframed}
    
    \item \textit{Optional:} Remember the n-dimensional  cube graph (See the lectures or Task 7). Find a Hamiltonian cycle in $Q_3$.  Does $Q_3$ have an Eulerian circuit?
   \begin{mdframed}\vspace{1 cm}\end{mdframed}
\end{enumerate}

\textbf{Task 2}

Two graphs are isomorphic if there exists a matching or relabelling of the vertices of the graphs.

\begin{center} \includegraphics[scale=.3]{Recitation 10/Iso_Ex.png}\end{center}

The two graphs above are isomorphic because we can match vertex a with 0, b with 3, c with 1, d with 4, and e with 2. Given this matching of the vertices, it's easy to check that the vertices are connected the same way in the graphs. Yet, it's hard to determine if two graphs are actually isomorphic! Let's try this ourselves:

\begin{enumerate}
    \item Out of the three graphs below, determine which two are isomorphic.
    
    Hint: Which of these are bipartite graphs?
    \begin{center} \includegraphics[scale=.44]{Recitation 10/isomorphism_prob.png}\end{center}
    \begin{mdframed}\vspace{1 cm}\end{mdframed}
    
    \item \textit{Optional:} Let's try an even harder example! Determine which two out of these four graphs are isomorphic.
    \begin{center} \includegraphics[scale=.22]{Recitation 10/Isomorphism_Problem.png}\end{center}
    \begin{mdframed}\vspace{1 cm}\end{mdframed}
\end{enumerate}

\textbf{Task 3}
\begin{enumerate}[a.]
    \item Given a certain labelling of the vertices, how many distinct graphs can be made with $4$ vertices?
    
    \begin{mdframed}\vspace{1 cm}\end{mdframed}
    
    \item How many graphs with 4 vertices exist if we don't count repetitions under isomorphism?
    \begin{mdframed}\vspace{1 cm}\end{mdframed}
    
    \item For any graph $G$, explain why $$\sum_{v \in V(G)}deg(v) = 2|E(G)|$$
    \begin{mdframed}\vspace{2 cm}\end{mdframed}
    
    \item The Cafeteria is having an event where two customers that purchase a sandwich to split receive 10\% off their order. The 35 CS22 crewmates are planning on taking advantage of this offer by having each crewmate split sandwiches with 3 other crewmates. Prove that this is impossible.
    \begin{mdframed}\vspace{2 cm}\end{mdframed}
    
    \item \textit{Optional:} Prove that in any graph there exist two vertices with the same degree.
    
    \textit{Hint:} What are the possible degrees of a vertex in a graph with $n$ nodes?
    \begin{mdframed}\vspace{2 cm}\end{mdframed}
\end{enumerate}

\textbf{Checkpoint 1:} Call over a TA.
\newpage

\section*{Trees!}

\begin{itemize}
    \item A \textbf{cycle} is a path that starts and ends with the same vertex. 
    %A cycle is \textbf{simple} if there are no other repeated vertices.
    \item A graph is \textbf{cyclic} if it contains a simple cycle and \textbf{acyclic} otherwise.
    \item A graph is \textbf{connected} if there is a simple path between each pair of vertices (that is, all vertices are connected).
    \item A \textbf{tree} is a connected, acyclic graph.
    \item A \textbf{forest} is an acyclic graph. In other words, a forest is a set of trees.
    \item A \textbf{leaf} of a tree is a vertex with degree 1.
    \item A \textbf{spanning tree} of a graph $G$ is a subgraph $T$ of $G$ such that $V(T) = V(G)$ and $T$ is a tree.
\end{itemize}

\textbf{Task 4}

In 1997, a chess-playing computer program called Deep Blue beat the reigning world champion of chess, Garry Kasparov, in a 6-round series. Although it was unclear whether computers can reason as well as humans, the victory did establish that computers are better at evaluating complex game trees, the basis of computer-based game strategies.

A game tree shows all possible playing strategies of both players in a game. Each vertex in the tree represents a configuration of the game as well as an indication of whose turn it is to play next. We illustrate game trees in a much simpler context of tic-tac-toe.

\begin{center}\caption{The first three layers of the game tree of tic-tac-toe} \includegraphics[scale=.3]{Recitation 10/tictactoe_init.png}\end{center}
The \textbf{root} of the tree is the initial configuration of the game, which is an empty grid in the case of tic-tac-toe. The \textbf{children} of a configuration c are all the configurations that can be reached from c by a single move of the correct player. A configuration is a \textbf{leaf} vertex in the tree if the game is over.

Here, we see a subtree near the end of the game:

\begin{center}\includegraphics[scale=.3]{Recitation 10/tictactoe_tree.png}\end{center}

\begin{enumerate}
    \item How many leaf nodes are in the subtree drawn above?
    \begin{mdframed}\vspace{1 cm}\end{mdframed}
    \item What is the length of the longest path from a root to a leaf in a tree represented by the full tic-tac-toe game?
    \begin{mdframed}\vspace{1 cm}\end{mdframed}

    \item \textit{Optional:} What is the length of the longest path between any two vertices in a tree represented by the full tic-tac-toe game?
    \begin{mdframed}\vspace{1 cm}\end{mdframed}
\end{enumerate}

\textbf{Task 5}
\begin{enumerate}
    \item Prove that in a tree, there is exactly one simple path between any two vertices.
    \begin{mdframed}\vspace{3 cm}\end{mdframed}
    
    \item One way of defining a tree is as a \textbf{minimally connected, maximally acyclic graph}. This means that if we remove an edge, our graph will no longer be connected and if we add an edge, the graph will no longer be acyclic. Prove these two statements.
    \begin{mdframed}\vspace{2 cm}\end{mdframed}

    \item \textit{Optional:} Let $T$ be a tree with at least two vertices. Prove that $T$ has at least two leaves.
    \textit{Hint:} Start by assuming that the longest simple path in $T$ does not start and end with leaves. What contradiction can we reach?
    \begin{mdframed}\vspace{2 cm}\end{mdframed}
\end{enumerate}

\section*{Induction on Graphs}

Many graph theory questions involve induction. When doing induction on graphs, we often use a technique referred to in CS22 as "build-down" induction. This is just regular induction, with a slightly different perspective. To review, here are the steps of an inductive proof:

\begin{enumerate}
    \item Define the predicate $P(n)$.

    \item Show that the base case is true.

    \item Assume the inductive hypothesis is true. If you are using standard induction then you will assume 
		$P(k)$ is true for some integer $k$. If you are using strong induction then you will assume $P(i)$ is true for all $i \leq k$.

    \item Show that $P(k+1)$ is true given the inductive hypothesis.
    
        Here is where build down induction differs from our usual perspective. A lot of the time, we start by invoking our inductive hypothesis, and manipulating the $k$ case to get the $k+1$ case. In build down induction, we do the following instead:

        \begin{enumerate}[i.]
            \item Consider an arbitrary $k+1$ case.
            \item Alter the $k+1$ case to obtain an $k$ case.
            \item Invoke the induction hypothesis, demonstrating that your property holds for the $k$ case.
            \item Recover the original $k+1$ case by undoing your alteration.
            \item Prove that the property still holds despite building back up to the original $k+1$ case.
        \end{enumerate}
\end{enumerate}

This technique is especially useful when it is not clear how to address all $k+1$ cases in general when starting from an $k$ case. For example, if we started with a graph with $k$ edges, where do we add an edge to obtain a general graph with $k+1$ edges? There are a lot of ways to add this extra edge that we have to consider. To avoid this, we can start with an arbitrary $k+1$ case. Then, we build-down to an $k$ case, at which point we can invoke our induction hypothesis. Then, we build back up to the \emph{original} $k+1$ case.

Try it out on a few sample problems!

\textbf{Task 6}
\begin{enumerate}
    \item Prove that a tree with $n$ vertices has $n-1$ edges.
    \begin{mdframed}\vspace{4 cm}\end{mdframed}

    \item \textit{Optional:} Prove that $K_n$ has $\frac{n(n-1)}{2}$ edges. \textit{Note that we could do this without build-down induction, but you should use it here for practice.}
   \begin{mdframed}\vspace{4 cm}\end{mdframed}
\end{enumerate}

\newpage
\section*{\textit{Optional:} Graph Coloring}
\textit{Note: This entire section is optional, so feel free to get checked off now. Graph coloring will \textbf{not be on the final}, but we think it's a super interesting topic in graph theory!}

Maps are often drawn such that no two adjacent regions are the same color. This makes it easy when looking at the map to see the distinct regions. For instance, here is a map of the United States where no adjacent states are the same color:

\begin{center}
\includegraphics[scale=.25]{map-us.png}
\end{center}

You might notice that this map was drawn with only 4 distinct colors. There are 50 U.S. states, so it might be a bit surprising that this is possible. This leads to a question that cartographers and mathematicians have asked for centuries: given a map, how many colors will we need to draw it?

Like so much else in life, this is really a graph theory problem in disguise! Any map can be easily converted to a graph.

For instance, we can represent this map of a region of Europe with the displayed graph. The graph was created by making each country into a vertex, and drawing an edge between two vertices if the countries share a border.

\begin{center}
\includegraphics[scale=.25]{map-europe.png}
\hspace{1cm}
\includegraphics[scale=.4]{graph-europe.PNG}
\end{center}
\newpage
\textbf{Task 7}
\begin{enumerate}
    \item Recall the n-dimensional cube graph; in particular, consider $Q_1, Q_2$, and $Q_3$, shown below:
    \begin{center} \includegraphics[scale=.5]{cube_graph.png} \end{center}
    What is the minimum number of colors we need to properly color $Q_1$, $Q_2$, and $Q_3$? What about for $Q_n$?
    \begin{mdframed}\vspace{1 cm}\end{mdframed}
    
    \item In task 2, we determined which two out of these three graphs are isomorphic:
    \begin{center}\includegraphics[scale=.5]{Recitation 10/isomorphism_prob.png}\end{center}
    Graph coloring is often helpful in determining isomorphic graphs. Determine the minimum number of colors needed to properly color $G_1$, $G_2$, and $G_3$.
    \begin{mdframed}\vspace{1 cm}\end{mdframed}
\end{enumerate}

\textbf{Checkoff: Call over a TA!}

\end{document}
